% GENERATED FILE. Edit canonical lessons, then run npm run generate:books.
% canonical-source-hash: fnv1a64:2cf621e72a7506d8
% canonical-lessons: GE-C361-einwanderung, GE-C361-flucht, GE-C361-fluechtling, GE-C361-integrieren, GE-C361-gewalt, GE-R361-review-354-361

\chapter{Ein neues Land}
\label{ch:ge-ein-neues-land}
\hlchaptermodality{361}

\begin{chapteropening}
\textbf{By the end of this chapter:} \emph{I can talk about immigration, flight and refugees, integration and violence.}

Complete the last lesson of chapter 361: I can talk about immigration, flight and refugees, integration and violence.
\end{chapteropening}

\section[die Einwanderung]{die Einwanderung — immigration}

\label{lesson:GE-C361-einwanderung}



\begin{quote}
Before the new one: say the German for to discriminate against, then the German for equal rights.
\end{quote}

\subsection*{You'll want to know: die Einwanderung}
\textbf{die Einwanderung} — \textquotedblleft{}immigration\textquotedblright{}.

\begin{quote}
Die Stadt ist durch Einwanderung schnell gewachsen.
\end{quote}

\begin{quote}
\emph{The city has grown quickly through immigration.}
\end{quote}

From \emph{einwandern} — to immigrate, with \emph{sein}. \emph{Der Einwanderer} — the immigrant.

\begin{cousinweb}
\emph{Einwanderung} is the noun of \emph{einwandern} \textquotedblleft{}to immigrate\textquotedblright{}, built on \emph{wandern} \textquotedblleft{}to wander\textquotedblright{}, inherited from Old High German \emph{wandarōn} and Proto-Germanic \emph{wandrōną} (EtymDB).
\end{cousinweb}

\begin{grammarlens}[title={What you've built}]
One more word for telling, explaining and describing at B1.
\end{grammarlens}

\subsection*{Your turn}
\begin{itemize}
\raggedright
  \item \emph{Say it:} \emph{die Einwanderung}
  \item \emph{Say it:} \emph{die Einwanderung}, once more
  \item \emph{Say it:} a sentence of your own with \emph{die Einwanderung}
  \item \emph{Recall:} say the German for to discriminate against, then the German for equal rights, then say \emph{die Einwanderung} again
\end{itemize}

\begin{tcolorbox}[breakable,colback=teal!4,colframe=teal!35!black,title={Before you move on}]
What does die Einwanderung mean? (\textquotedblleft{}Immigration\textquotedblright{}.) Say it once more.
\end{tcolorbox}
\section[die Flucht]{die Flucht — flight, escape}

\label{lesson:GE-C361-flucht}



\begin{quote}
Before the new one: say the German for equal rights, then the German for immigration.
\end{quote}

\subsection*{You'll want to know: die Flucht}
\textbf{die Flucht} — \textquotedblleft{}flight, escape\textquotedblright{}.

\begin{quote}
Nach der Flucht lebte die Familie zwei Jahre in einem Lager.
\end{quote}

\begin{quote}
\emph{After their flight the family lived in a camp for two years.}
\end{quote}

\emph{Auf der Flucht} — on the run. The verb is \emph{fliehen} — to flee.

\begin{cousinweb}
\emph{Flucht} is inherited from Middle High German \emph{vluht}, Old High German \emph{fluht} and Proto-Germanic \emph{fluhtiz} \textquotedblleft{}flight\textquotedblright{} (EtymDB), the noun of the verb for fleeing.
\end{cousinweb}

\begin{grammarlens}[title={What you've built}]
One more word for telling, explaining and describing at B1.
\end{grammarlens}

\subsection*{Your turn}
\begin{itemize}
\raggedright
  \item \emph{Say it:} \emph{die Flucht}
  \item \emph{Say it:} \emph{die Flucht}, once more
  \item \emph{Say it:} a sentence of your own with \emph{die Flucht}
  \item \emph{Recall:} say the German for equal rights, then the German for immigration, then say \emph{die Flucht} again
\end{itemize}

\begin{tcolorbox}[breakable,colback=teal!4,colframe=teal!35!black,title={Before you move on}]
What does die Flucht mean? (\textquotedblleft{}Flight, escape\textquotedblright{}.) Say it once more.
\end{tcolorbox}
\section[der Flüchtling]{der Flüchtling — the refugee}

\label{lesson:GE-C361-fluechtling}



\begin{quote}
Before the new one: say the German for immigration, then the German for flight.
\end{quote}

\subsection*{You'll want to know: der Flüchtling}
\textbf{der Flüchtling} — \textquotedblleft{}the refugee\textquotedblright{}.

\begin{quote}
Die Schule hat zwanzig junge Flüchtlinge aufgenommen.
\end{quote}

\begin{quote}
\emph{The school has taken in twenty young refugees.}
\end{quote}

Plural: \emph{die Flüchtlinge}. Today people also say \emph{der Geflüchtete, die Geflüchtete}.

\begin{cousinweb}
\emph{Flüchtling} is built on \emph{flüchten} \textquotedblleft{}to flee\textquotedblright{}, from \emph{Flucht} \textquotedblleft{}flight\textquotedblright{}, Old High German \emph{fluht} (EtymDB), with the ending \emph{-ling} for a person.
\end{cousinweb}

\begin{grammarlens}[title={What you've built}]
One more word for telling, explaining and describing at B1.
\end{grammarlens}

\subsection*{Your turn}
\begin{itemize}
\raggedright
  \item \emph{Say it:} \emph{der Flüchtling}
  \item \emph{Say it:} \emph{der Flüchtling}, once more
  \item \emph{Say it:} a sentence of your own with \emph{der Flüchtling}
  \item \emph{Recall:} say the German for immigration, then the German for flight, then say \emph{der Flüchtling} again
\end{itemize}

\begin{tcolorbox}[breakable,colback=teal!4,colframe=teal!35!black,title={Before you move on}]
What does der Flüchtling mean? (\textquotedblleft{}The refugee\textquotedblright{}.) Say it once more.
\end{tcolorbox}
\section[integrieren]{integrieren — to integrate}

\label{lesson:GE-C361-integrieren}



\begin{quote}
Before the new one: say the German for flight, then the German for the refugee.
\end{quote}

\subsection*{You'll want to know: integrieren}
\textbf{integrieren} — \textquotedblleft{}to integrate\textquotedblright{}.

\begin{quote}
Der Verein hilft, neue Nachbarn zu integrieren.
\end{quote}

\begin{quote}
\emph{The club helps to integrate new neighbours.}
\end{quote}

\emph{Sich integrieren} — to fit in. Perfekt: \emph{hat integriert}.

\begin{cousinweb}
\emph{integrieren} is a learned borrowing from Latin \emph{integrāre} \textquotedblleft{}to make whole, to renew\textquotedblright{} (EtymDB; Lewis \& Short), from \emph{integer} \textquotedblleft{}whole, untouched\textquotedblright{}.
\end{cousinweb}

\begin{grammarlens}[title={What you've built}]
One more word for telling, explaining and describing at B1.
\end{grammarlens}

\subsection*{Your turn}
\begin{itemize}
\raggedright
  \item \emph{Say it:} \emph{integrieren}
  \item \emph{Say it:} \emph{integrieren}, once more
  \item \emph{Say it:} a sentence of your own with \emph{integrieren}
  \item \emph{Recall:} say the German for flight, then the German for the refugee, then say \emph{integrieren} again
\end{itemize}

\begin{tcolorbox}[breakable,colback=teal!4,colframe=teal!35!black,title={Before you move on}]
What does integrieren mean? (\textquotedblleft{}To integrate\textquotedblright{}.) Say it once more.
\end{tcolorbox}
\section[die Gewalt]{die Gewalt — violence; power, force}

\label{lesson:GE-C361-gewalt}



\begin{quote}
Before the new one: say the German for the refugee, then the German for to integrate.
\end{quote}

\subsection*{You'll want to know: die Gewalt}
\textbf{die Gewalt} — \textquotedblleft{}violence; power, force\textquotedblright{}.

\begin{quote}
Gewalt ist keine Lösung.
\end{quote}

\begin{quote}
\emph{Violence is no solution.}
\end{quote}

No plural in this sense. Also power: \emph{die Gewalt über das Auto verlieren} — to lose control of the car.

\begin{cousinweb}
\emph{Gewalt} is inherited from Middle High German \emph{gewalt} \textquotedblleft{}strength, power\textquotedblright{} and Old High German \emph{giwalt} (EtymDB), from Proto-Germanic \emph{waldą}, built on \emph{waldaną} \textquotedblleft{}to rule\textquotedblright{}.
\end{cousinweb}

\begin{grammarlens}[title={What you've built}]
One more word for telling, explaining and describing at B1.
\end{grammarlens}

\subsection*{Your turn}
\begin{itemize}
\raggedright
  \item \emph{Say it:} \emph{die Gewalt}
  \item \emph{Say it:} \emph{die Gewalt}, once more
  \item \emph{Say it:} a sentence of your own with \emph{die Gewalt}
  \item \emph{Recall:} say the German for the refugee, then the German for to integrate, then say \emph{die Gewalt} again
\end{itemize}

\begin{tcolorbox}[breakable,colback=teal!4,colframe=teal!35!black,title={Before you move on}]
What does die Gewalt mean? (\textquotedblleft{}Violence; power, force\textquotedblright{}.) Say it once more.
\end{tcolorbox}
\section[(dialogue)]{Words from chapters 354–361 — a review}

\label{lesson:GE-R361-review-354-361}



\begin{quote}
Say the German for the meeting and the session.
\end{quote}

\subsection*{Your turn}
Say these aloud:

\begin{itemize}
\raggedright
  \item the meeting, the session, the minutes, to present, the presentation
  \item the project, the order, to work together, the cooperation, to organize
  \item the career, the rise, the challenge, to cope with, to fail
  \item the trade union, the works council, the demand, to support, the support
  \item society, democracy, to vote, the vote, to contribute
  \item justice, inequality, poverty, wealth, to put at a disadvantage
  \item the prejudice, tolerance, the minority, to discriminate against, equal rights
  \item immigration, flight, the refugee, to integrate, violence
\end{itemize}

\begin{tcolorbox}[breakable,colback=teal!4,colframe=teal!35!black,title={Before you move on}]
The meeting? (\textbf{die Besprechung}.) The session? (\textbf{die Sitzung}.) The minutes? (\textbf{das Protokoll}.) To present? (\textbf{präsentieren}.) The presentation? (\textbf{die Präsentation}.) The project? (\textbf{das Projekt}.) The order? (\textbf{der Auftrag}.) To work together? (\textbf{zusammenarbeiten}.) The cooperation? (\textbf{die Zusammenarbeit}.) To organize? (\textbf{organisieren}.) The career? (\textbf{die Karriere}.) The rise? (\textbf{der Aufstieg}.) The challenge? (\textbf{die Herausforderung}.) To cope with? (\textbf{bewältigen}.) To fail? (\textbf{scheitern}.) The trade union? (\textbf{die Gewerkschaft}.) The works council? (\textbf{der Betriebsrat}.) The demand? (\textbf{die Forderung}.) To support? (\textbf{unterstützen}.) The support? (\textbf{die Unterstützung}.) Society? (\textbf{die Gesellschaft}.) Democracy? (\textbf{die Demokratie}.) To vote? (\textbf{abstimmen}.) The vote? (\textbf{die Abstimmung}.) To contribute? (\textbf{beitragen}.) Justice? (\textbf{die Gerechtigkeit}.) Inequality? (\textbf{die Ungleichheit}.) Poverty? (\textbf{die Armut}.) Wealth? (\textbf{der Reichtum}.) To put at a disadvantage? (\textbf{benachteiligen}.) The prejudice? (\textbf{das Vorurteil}.) Tolerance? (\textbf{die Toleranz}.) The minority? (\textbf{die Minderheit}.) To discriminate against? (\textbf{diskriminieren}.) Equal rights? (\textbf{die Gleichberechtigung}.) Immigration? (\textbf{die Einwanderung}.) Flight? (\textbf{die Flucht}.) The refugee? (\textbf{der Flüchtling}.) To integrate? (\textbf{integrieren}.) Violence? (\textbf{die Gewalt}.)
\end{tcolorbox}
